% ECONOMICS_PROBLEM: {"id": "F32-247", "title": "Sudden-stop amplification", "jel_code": "F32", "source_id": "OP-0336"}
% !TeX program = xelatex
\documentclass[11pt,a4paper]{article}
\usepackage[margin=23mm]{geometry}
\usepackage{fontspec}
\setmainfont{Latin Modern Roman}
\usepackage{amsmath,amssymb,mathtools}
\usepackage{xcolor,enumitem,booktabs,longtable,array,xurl,hyperref}
\definecolor{muted}{HTML}{53636C}
\hypersetup{colorlinks=true,urlcolor=blue,linkcolor=blue}
\setlength{\parindent}{0pt}
\setlength{\parskip}{5pt}
\newcommand{\R}{\mathbb R}
\newcommand{\E}{\mathbb E}
\newcommand{\N}{\mathbb N}
\newcommand{\ind}{\mathbf 1}
\newcommand{\Var}{\operatorname{Var}}
\newcommand{\Cov}{\operatorname{Cov}}
\newcommand{\supp}{\operatorname{supp}}
\DeclareMathOperator*{\argmin}{arg\,min}
\DeclareMathOperator*{\argmax}{arg\,max}
\newcommand{\entrymeta}[3]{{\sffamily\small\color{muted}\textbf{\texttt{#1}}\quad|\quad #2\par #3\par}}
\newcommand{\Problemref}[1]{\texttt{#1}}
\newcommand{\JELref}[2]{\texttt{#2} (JEL \texttt{#1})}
\begin{document}
\section*{Sudden-stop amplification}
\textbf{Problem ID:} \texttt{F32-247}\quad \textbf{JEL:} \texttt{F32}\par
% BEGIN ECONOMICS STATEMENT
\label{problem:OP-0336}
{\sffamily\small\color{muted}Original subject group: International finance and exchange rates\par}
\label{definitions:problem:30007618}
\textbf{Audit classification:} Background; proposed extension. The primary paper already derives model-specific network amplification and reports calibrated crisis effects. The broader external-financing intervention is editorial; its relaxed-constraint denominator may vanish.
\subsection*{Definitions and precise research target}
Fix a small open economy with firms linked by an input-output matrix \(A\), external debt denominated in stated currencies, banks and collateral constraints. A sudden-stop intervention reduces external financing supply by \(u>0\) at date \(t\), independently of specified domestic productivity news. Let \(Y_h(u;A,b)\) be output at horizon \(h\) under this intervention, where \(b\) collects initial maturity, currency mismatch and collateral exposures. Define amplification relative to a counterfactual economy with relaxed financing constraints and the same technologies by
\[\Gamma_h(A,b)=\frac{E[Y_h(0;A,b)-Y_h(u;A,b)]}{E[Y_h^0(0;A,b)-Y_h^0(u;A,b)]},\]
when the denominator is positive and the relaxed-constraint economy retains a stated transmission channel; otherwise report the signs and the loss difference defined below. Superscript \(0\) denotes that relaxed-constraint economy. The task is to identify amplification and characterize network features increasing crisis losses. Model endogenous exchange rates, collateral prices, substitution and lender risk-bearing; holding these fixed changes the target. Establish the financing shock's exclusion and avoid conflating capital withdrawal with investors responding to output declines. Report nonlinearities as constraints become binding, sensitivity to input substitutability and possible firm exit. Where only total responses are identified, distinguish them from model-dependent counterfactual amplification. An input-output correlation with crisis severity alone does not isolate a causal network mechanism.

\textbf{Additional formal definitions and scope (editorial).} Define \(A_{ij}\) as the input share from sector \(i\) in sector \(j\), with nonnegative entries and feasibility conditions from the chosen production functions; Cobb--Douglas share restrictions and CES substitutability cannot be mixed implicitly. The sudden-stop intervention changes a named financing-supply limit or price equation. Relaxing collateral constraints can remove that shock's entire effect, so the denominator of \(\Gamma_h\) may be zero. Use the ratio only when the frictionless benchmark retains a nonzero financing-price or expenditure channel, and specify that channel. Otherwise report the difference in losses \(\mathcal A_h=E[Y_h(0)-Y_h(u)]-E[Y_h^0(0)-Y_h^0(u)]\). This is a model-specific amplification contrast, not an automatically identified causal network effect. Define the same exogenous shock, technology and initial resource state for both comparisons.

For the identification part, declare an admissible class \(\Theta\) of structural models, including preferences, technologies, shock laws, measurement/sampling rules and any equilibrium selector. If \(O\) denotes the complete observable history and \(T(\theta)\) the target defined above, the identified set is \(\mathcal I_T=\{T(\theta):\theta\in\Theta,\ \mathcal L_\theta(O)=\mathcal L_{\rm obs}(O)\}\); point identification means this set is a singleton. A \(do(a)\) comparison replaces stated policy/technology equations while fixing the declared exogenous law and re-solving the remaining equations. These definitions do not assert that an identification theorem holds without further restrictions.
\subsection*{Variants and assumption changes}
\begin{enumerate}
\item \textbf{Two-sector collateral model (established source subcase).} Characterize the source documented network-dependent amplification within its stated tradable/nontradable model; its theoretical and calibrated answers are already available.
\item \textbf{External financing-supply intervention (editorial).} Extend to many sectors and currency mismatch; use loss differences when removing constraints makes the frictionless shock response zero.
\end{enumerate}
\subsection*{References and source audit}
\begin{enumerate}
\item NBER indexed record verifies four authors, 2025 WP 34604 and DOI. Boston Fed WP 26-1 (2026), December 2025 text, inspected with its distinct DOI 10.29412/res.wp.2026.01. Locator: Boston Fed version Abstract and Sections 2--5. Support: Model-specific network amplification results. NBER direct fetch failed; equivalent primary Boston Fed PDF accessible. URL: \url{https://www.nber.org/papers/w34604}.
\item NBER indexed record verifies four authors, 2025 WP 34604 and DOI. Boston Fed WP 26-1 (2026), December 2025 text, inspected with its distinct DOI 10.29412/res.wp.2026.01. Locator: Boston Fed version Abstract and Sections 2--5. Support: Model-specific network amplification results. NBER direct fetch failed; equivalent primary Boston Fed PDF accessible. URL: \url{https://www.bostonfed.org/-/media/Documents/Workingpapers/PDF/2026/WP2601.pdf}.
\end{enumerate}
\begin{enumerate}\raggedright
\item\hypertarget{definitions:source:30007618:1}{} Jorge Miranda-Pinto; Eugenio Rojas; Felipe Saffie; Álvaro Silva. 2025. \emph{Connected for Better or Worse? The Role of Production Networks in Financial Crises}. Boston Fed version Abstract and Sections 2--5.
\url{https://www.nber.org/papers/w34604}.
DOI: \url{https://doi.org/10.3386/w34604}.
\item\hypertarget{definitions:source:30007618:2}{} Jorge Miranda-Pinto; Eugenio Rojas; Felipe Saffie; Álvaro Silva. 2026. \emph{Connected for Better or Worse? The Role of Production Networks in Financial Crises}. Boston Fed version Abstract and Sections 2--5.
\url{https://www.bostonfed.org/-/media/Documents/Workingpapers/PDF/2026/WP2601.pdf}.
DOI: \url{https://doi.org/10.29412/res.wp.2026.01}.
\end{enumerate}

\subsection*{Common conventions: Source mathematical conventions}

These conventions apply to each entry unless explicitly replaced.

\textbf{Models and identification.} A primitive model \(m\) specifies all probability laws, feasible choices, information, equilibrium and measurement rules used by its target. The admissible class \(\mathcal M\) is fixed independently of outcomes. If \(P_O(m)\) is its observable probability law and \(\tau(m)\) the target, its identified set at observable law \(P\) is
\[
 \mathcal I_\tau(P)=\{\tau(m):m\in\mathcal M,\ P_O(m)=P\}.
\]
Point identification means that this set is a singleton. An empty set signals incompatibility of data and assumptions. Coordinate bounds need not describe the joint set. Bounds are sharp only if every retained target is attainable or arbitrarily approachable by compatible models. Finite bounds require explicit restrictions on unobserved quantities; unrestricted confounding, hidden wealth, extrapolation or arbitrary equilibrium selection can yield uninformative sets.

\textbf{Causal interventions.} A potential outcome \(Y_i(p)\) is the outcome under a complete policy regime \(p\), including timing, financing, information and market spillovers. The target population and units are fixed before treatment. With interference, use the complete assignment vector or a justified exposure mapping; individual no-interference notation alone cannot identify an economy-wide intervention. A difference of mean potential outcomes does not require the cross-world joint distribution. Conditioning on post-treatment migration, survival, enrollment or employment changes the estimand and needs separate justification.

\textbf{Statistical statements.} An estimator, confidence set, forecast or test specifies a sampling law, model class and loss/coverage criterion. Uniform \((1-\alpha)\)-coverage means \(\inf_{m\in\mathcal M}\Pr_m\{\tau(m)\in C_n\}\geq1-\alpha\), or an explicitly asymptotic version. The whole parameter space trivially has coverage, so informativeness requires a stated width, risk or power objective. A forecasting improvement is relative to a named benchmark, information set and evaluation population.

\textbf{Equilibrium and optimization.} An equilibrium comprises feasible plans, optimality, beliefs where needed, budget constraints and all market/resource clearing. The primitives or a stated benchmark define each correspondence \(\mathcal E(p;m)\). Nonemptiness is assumed or proved, never inferred just from compact policy sets. A selected-equilibrium target uses a declared selection rule; a robust target quantifies over all permitted equilibria. Use suprema or infima unless attainment is established. An \(\varepsilon\)-optimal policy is feasible and within \(\varepsilon\) of the finite optimum value. Report an empty feasible set separately.

\textbf{Welfare and accounting.} Normative utility, interpersonal normalization, social weights, numeraire, discounting and the use of public receipts are primitives. Behavioral utility need not coincide with the welfare criterion. Household welfare includes ownership income and rents once; adding profits again to compensating variation generally double-counts them. A monetary equivalent is defined by an explicit transfer and price convention, with monotonicity and range conditions. Average effects, mechanism contributions, transfers and welfare losses are different targets. Infinite sums require convergence and dynamic feasible plans require terminal or no-Ponzi conditions.

\textbf{Source versions.} A working-paper date, revision date and journal publication year are distinguished. A DOI for a journal version is not automatically the DOI of an earlier manuscript. Locators refer to the stated version and printed pages where specified. A metadata-only check does not verify a claim about a full-text conclusion. Every entry's source subsection narrows the provenance claim accordingly.
% END ECONOMICS STATEMENT
\end{document}
